\centerline{\bf \Huge Равенства: работаем ручками.}


\problem{1}
 Для различных положительных действительных чисел $a, b$ справедливо равенство

$$
\frac{a}{a^3+a+1}=\frac{b}{b^3+b+1} .
$$


Найдите значение выражения

$$
\frac{13-a^2 b-b^2 a}{2+a^2 b+b^2 a}
$$

\problem{2}
Найти все пары цельх неотрицательных чисел $(k, m)$, являющихся решениями уравнения

$$
2 k^2+7 k=2 m k+3 m+36
$$

\problem{3}
Найдите все пары целых чисел $(x, y)$, удовлетворяющие уравнению

$$
\left(x^2+y^2\right)(x+y-3)=2 x y
$$

\problem{4}
 Действительные числа $x, y, z$ выбираются так, что выполняются равенства $x y+y z+z x=4, x y z=6$. Докажите, что при любом таком выборе значение выражения

$$
\left(x y-\frac{3}{2}(x+y)\right)\left(y z-\frac{3}{2}(y+z)\right)\left(z x-\frac{3}{2}(z+x)\right)
$$


является одним и тем же числом, и найти это число.
\problem{5}
Натуральные числа $x, y$ таковы, что $2 x^2+x=3 y^2+y$. Докажите, что число $x-y$ является точным квадратом.
\problem{6}
 Натуральные числа $x$ и $y$ таковы, что $3 x^2+3 x+1=y^2$. Докажите, что $y$ представляется в виде суммы квадратов двух последовательных натуральных чисел.